% Robuste matte-/spesialtegn: render alltid via mattefont, uansett tekstfont
\usepackage{newunicodechar}
\newunicodechar{≈}{\ensuremath{\approx}}
\newunicodechar{≠}{\ensuremath{\neq}}
\newunicodechar{≤}{\ensuremath{\leq}}
\newunicodechar{≥}{\ensuremath{\geq}}
\newunicodechar{∈}{\ensuremath{\in}}
\newunicodechar{μ}{\ensuremath{\mu}}
\newunicodechar{β}{\ensuremath{\beta}}
\newunicodechar{α}{\ensuremath{\alpha}}
\newunicodechar{γ}{\ensuremath{\gamma}}
\newunicodechar{τ}{\ensuremath{\tau}}
\newunicodechar{κ}{\ensuremath{\kappa}}
\newunicodechar{×}{\ensuremath{\times}}
\newunicodechar{−}{\ensuremath{-}}
\newunicodechar{→}{\ensuremath{\rightarrow}}
\newunicodechar{∑}{\ensuremath{\sum}}

% Forsidetall — oppdateres automatisk av _assets/bygg_rapport.py
\newcommand{\antallsider}{69}
\newcommand{\antallord}{14\,571}

% Tabeller: mindre skrift + litt mindre kolonneavstand, så brede tabeller får plass
\usepackage{etoolbox}
\AtBeginEnvironment{longtable}{\small}
\AtBeginEnvironment{tabular}{\small}
\setlength{\tabcolsep}{5pt}

% Naturlige sideskift: hvert hovedkapittel starter på ny side, og
% front-matter (egenerklæring, Sammendrag, Abstract, Innhold, Figur-/Tabelliste)
% får egne sider — men vanlige underkapitler (5.1, 6.1, ...) brytes ikke.
% Tittelsiden er en rå LaTeX-blokk (ikke en \section), så front-matter starter «på».
% Første ekte \section (1.0 Innledning) slår front-matter av.
\newcounter{secidx}
\newif\iffrontmatter\frontmattertrue
\let\oldsection\section
\renewcommand{\section}{%
  \stepcounter{secidx}%
  \ifnum\value{secidx}=1 \frontmatterfalse\fi
  \clearpage\oldsection}
\let\oldsubsection\subsection
\renewcommand{\subsection}{\iffrontmatter\clearpage\fi\oldsubsection}

% Figur-flyt: la figurene plasseres uten å etterlate store tomrom
\usepackage{float}
\renewcommand{\topfraction}{0.9}
\renewcommand{\bottomfraction}{0.9}
\renewcommand{\textfraction}{0.08}
\renewcommand{\floatpagefraction}{0.7}
